\section{Fundamentals of Unonstrained Optimization}
The concept of a local minimizer is introduced,
\begin{Def}[Weak local minimizer]
    A point $x^*$ is a \textit{weak local minimizer} if there is a neighborhood
    % ------------------- Footnote ---------------------
    \footnote{
        A neighborhood of a point $x$, sometimes denoted $\mathcal{N}_x$ is an open set containing $x$. 
        Defining which subsets of are considered open is the act of choosing a topology for your space. 
        In this book they use the topology induced on $\mathbb{R}^n$ by the euclidean metric $\left( \| \cdot\|_2\right) : x \mapsto \sqrt{x^Tx}$. 
        For more on Euclidean topology, see \href{https://en.wikipedia.org/wiki/Real_coordinate_space\#Topological_properties}{\textcolor{blue}{here}}.
} 
% -------------------------------------------------------
$\mathcal{N}$ of $x^*$ such that $f\left(x^*\right) \leq f(x), \quad\forall x\in\mathcal{N}$. 
\end{Def}
On the other hand a \textit{strict} local optimizer is where the inequality above is strict.
A local minimizer $x^*$ is called \textit{isolated} if there exists a neighborhood in which $x^*$ is the only local minimizer.
Sufficient conditions for a point to be a local minimizer were presented as the following 
\begin{Thm}[Second-Order Sufficient Conditions]
    Suppose that $\nabla^2f$ is continous in an open neighborhood of $x^*$ and that $\nabla f(x^*) = 0$ and $\nabla^2f(x^*)$ is positive definite. Then $x^*$ is a strict local minimizer of $f$
\end{Thm}
\subsection{Questions}
\begin{enumerate}[label = \roman*)]
\item \textit{Do we use any of these?} Vplot has some unconstrained, least-squares algorithm.
\item \textit{What have you discussed previously on the topic of smooth optimization?} Interior point methods for optimizing SIMM delta.
\item \textit{How is the process of inverting a matrix affected by terrible condition number?} Which algorithms are used for matrix inversion? Is there some bound on iterations dependent on the condition number?
\end{enumerate}
\subsection{Solutions}
The exercises chosen for this chapter are 2.1, 2.3, 2.10, 2.11 and 2.12.
\subsubsection{Exercise 2.1}
Compute i) the gradient and ii) the Hessian of the Rosenbrock
\footnote{
    There are higher dimensional \href{https://en.wikipedia.org/wiki/Rosenbrock_function\#Multidimensional_generalizations}{\textcolor{blue}{generalizations}} of the Rosenbrock function which are classic problems to benchmark performance on $n$-dimensional, non-convex probelms. 
} function
\[
f(x) = 100(x_2 - x_1^2)^2 + (1-x_1)^2
\]
and iii) show that $x^* = (1,1)$ is the only local minimizer.
\begin{enumerate}[label = \roman*)]
    \item The gradient of a function is the vector of all its partial derivatives \[
        \nabla f (x) =\left(\pd[][f][x_1], \pd[][f][x_2]\right)^T.
    \]
    Computing the two derivatives in $x_1$ and $x_2$ we get
    \begin{align*}
        \pd[][f][x_1] &= -400(x_2 - x_1^2)x_1 - 2(1-x_1)\\
        \pd[][f][x_2] &= 200(x_2-x_1^2)
    \end{align*}
    \item The Hessian of a function is the matrix of all it's partial derivatives
    \[
    \textbf{H}_f = \left(\nabla^2 f\right)_{ij} = \pd[2][f][x_i,x_j] 
    \]
    As a mnemonic, $\nabla = \left(\pd[][][x_1], \dots, \pd[][][x_2]\right)$ and $\nabla^2 = \nabla\nabla^T$.
    \begin{align*}
        \pd[2][f][x_1] &= -400\left(x_2 - 3x_1^3\right) + 2\\
        \pd[2][f][x_1,x_2] = \pd[2][f][x_2,x_1] &= -400x_1\\
        \pd[2][f][x_2] &= 200 
    \end{align*}
    \item For a local minimizer, $\nabla f = 0$ and $\nabla^2f$ is positive definite.
    The first condition means to solve for which $x_1, x_2$ the equations
    \[
    \pd[][f][x_1] = 0, \quad \pd[][f][x_2] = 0
    \]
    hold.
    The second one is the easiest to start with, it implies the relation $x_2 = x_1^2$ for stationary points (a point for which $\nabla f(x) = 0$ is true) of $f$.
    Substituting into the former of the equations above we get
    \[
    2(1 - x_1) = 0 \implies x_1 = 1
    \]
    which in turn means $x_2 = 1$.
    This means there is only one stationary point and thus, at most one local minimizer.
    Only if the Hessian is positive definite can the stationary point be a local minimizer. Let's check!
    \[
    \nabla^2f(1,1) = 
    \begin{pmatrix}
        802 & -400\\
        -400 & 200
    \end{pmatrix}
    \]
    A symmetric matrix is positive definite if all it's leading principal minors are positive\footnote{S] \href{https://en.wikipedia.org/wiki/Sylvester\%27s_criterion}{\textcolor{blue}{Sylvester's criterion}}}. The principal minors are the determinants of square sub-matrices starting from the upper left corner. 
    The first one is the 1 \by 1 matrix in the top left corner, i.e. $802>0$.
    The second one is 
    \[
    \operatorname{det}\nabla^2 f(1,1) = 802\cdot200 - 400^2 = 400 > 0
    \]
    \hfill\sq
\end{enumerate}
\subsubsection{Exercise 2.3}
Starting with $f_1$,
\[
    \pd[][f_1][x_i]  = \pd[][a^Tx][x_i] = \pd[][][x_i]\sum_i a_ix_i = a_i.
\] 
So $\nabla f_1 = a\in\mathbb{R}^n$. 
Since the gradient is constant and the Hessian consists of second derivatives, it is zero $\nabla^2f_1 = \mathbf0$.

For $f_2$ we have
\begin{align*}
\pd[][f_2][x_k]  &= \sum_{i,j} A_{ij}\pd[][(x_ix_j)][x_k]  \\ &=\sum_{i,j}A_{ij}(\delta_{ik}x_j + x_i\delta_{jk}) \\
&=\sum_{j}A_{kj}x_j + \sum_i A_{ik}x_i \\
&=2A_kx
\end{align*}
where $A_k$ is the $k^{\text{th}}$ column of $A$ and $\delta_{ij}$ is the Kronecker delta\footnote{S] \href{https://en.wikipedia.org/wiki/Kronecker_delta}{\textcolor{blue}{Kronecker delta}}}. 
The last inequality follows from $A$ being symmetric.
Thus
\[
\nabla f_2 = (2A_1x,\dots,2A_nx)^T = 2 A x.
\]
For the Hessian, we need
\[
\pd[2][f_2][x_lx_k] = \pd[][][x_l]\left(2\sum_iA_{ik}x_i\right) = 2A_{lk}.
\]
which gives
\[
\nabla^2f_2 = 2A
\]
\subsubsection{Exercise 2.10}
\[
\pd[][\tilde{f}][z_i] = \sum_k\pd[][f][x_k]\pd[][x_k][z_i]
\]
To begin with, we compute the second factor in the sum above
\[
\pd[][x_k][z_i] = \pd[][][z_i]\left( \sum_jS_{kj}z_j + s_k\right) = S_{ki}.
\]
Substituting this for $\pd[][x_k][z_i]$ in the first equation we get
\[
\pd[][\tilde{f}][z_i] = \sum_k \pd[][f][x_k] S_{ki} = S_i^T \cdot \nabla f
\]
with $S_i$ being the $i^{\text{th}}$ column of $S$.
This means\footnote{If $S_i$ had b]n rows, we would have had $\nabla \tilde{f}(z) = S\nabla f(x)$}
\[
\nabla \tilde{f}(z) = \begin{bmatrix}
    S_1 \nabla f\\
    \vdots\\
    S_n \nabla f
\end{bmatrix} = S^T\nabla f(x)
\]
For the Hessian, we need
\[
\pd[2][f][z_i,z_j] = \sum_{k,l}\pd[2][f][x_k,x_l]\pd[][x_k][z_i]\pd[][x_l][z_j] = \sum_{k,l} \pd[2][f][x_k,x_l]S_{ki}S_{lj}
\]
Note that we could have hade second derivatives of $x$ w.r.t. $z$, but thankfully the first derivative is constant so we are off the hook!
The expression above can be rewritten as:
\[
\pd[2][f][z_i,z_j] = S_i^T \cdot \nabla^2 f \cdot  S_j
\]
So when assembling the whole Hessian matrix from the partial derivatives, we obtain
\[
\nabla \tilde{f}(z) = S^T \nabla f(x) S
\]
\hfill\sq 
\subsubsection{Exercise 2.11}
Assuming a unit step size, the relation between subsequent steps is 
\[
x_{k+1} = x_k + p
\]
We use this relation to make sure the affine relation between the iterates is preserved using this update rule.
To begin with, the relation between the iterates $z_k$ and $x_k$ is
\[
x_k = Sz_k + s.
\]
The step direction at each iteration is computed as
\[
p_k = -B_k^{-1}\nabla f(x_k)
\]
and 
\[
\tilde{p}_k = -\tilde{B}_k^{-1}\nabla \tilde{f}(z_k)
\]
Given that the initial Hessians approximations are chosen s.t $\tilde{B}_0 = S^TB_0S$ we can s] that
\[
\tilde{p}_0 =  -(S^TB_0S)^{-1}S^T\nabla f(x_0) = -S^{-1}B_0^{-1}\nabla f(x_0)
\]
which, for $z_{k+1}$, means
\begin{align*}
z_{k+1} =& z_k + \tilde{p_k} = z_k - S^{-1}B_k^{-1}\nabla f(x_k)\\
=& (S^{-1}(x_k - s)) - S^{-1}B_k^{-1}\nabla f(x_k)
\end{align*}
Applying the affine transformation $v \mapsto Sv +s$ on both sides
\[
\underbrace{x_k + B_{k}^{-1}\nabla f(x_k)}_{x_{k+1}} = Sz_{k+1} + s
\]
So, given $\tilde{B}_k = S^TB_kS$, the above holds. 
Since we can choose the initial approximations such that this holds, it remains to verify that the rank one and BFGS update rules preserve this relationship.
We begin with the rank-one update. 
\[
B_{k+1} \;=\; B_k \;+\; 
\frac{(y_k - B_k s_k)(y_k - B_k s_k)^{T}}
     {(y_k - B_k s_k)^{T} s_k}
\]
where
\[
s_k \;=\; x_{k+1} - x_k, 
\qquad 
y_k \;=\; \nabla f(x_{k+1}) - \nabla f(x_k).
\]
We know that for $k=0$, we have $x_k = Sz_k +s$ and $\tilde{B}_0 = S^TB_0S$.
If this update preserves the relation between the Hessian iterates, then the relation between the solution iterates will be preserved and continue to hold through further iterations.
The quantities $y_k$ and $s_k$ relate to their $\tilde{}$ counterparts as follows:
\[
\tilde{s}_k \;=\; z_{k+1} - z_k = S^{-1}(x_{k+1} -x_k) = S^{-1}s_k, 
\qquad 
\tilde{y}_k \;=\; \nabla f(z_{k+1}) - \nabla f(z_k) = S^Ty_k.
\]
We now show the update satisfies behaves desirably by substituting in the above relations,
\begin{align*}
    \tilde{B}_{k+1} &= \tilde{B}_k + \frac{(\tilde{y}_k - \tilde{B}_k \tilde{s}_k)(\tilde{y}_k - \tilde{B}_k \tilde{s}_k)^{T}}
     {(\tilde{y}_k - \tilde{B}_k \tilde{s}_k)^{T} \tilde{s}_k}\\
     &= S^TB_kS + \frac{(S^Ty_k - S^T B_k S S^{-1} s_k)(S^Ty_k - S^TB_kS S^{-1} s_k)^{T}}
     {(S^Ty_k - S^TB_kSS^{-1} s_k)^{T} S^{-1} s_k}\\
     &= S^TB_kS + \frac{S^T(y_k -  B_k s_k)(y_k - B_ks_k)^{T}S}
     {(y_k - B_k s_k)^{T} S S^{-1} s_k} \\
     &= S^T\underbrace{\left(B_k + \frac{(y_k -  B_k s_k)(y_k - B_ks_k)^{T}}
     {(y_k - B_k s_k)^{T}  s_k}\right)}_{B_{k+1}}S
\end{align*}
The procedure for the BFGS update rule is identical.\hfill\sq
\subsubsection{Exercise 2.12}
This question is not entirely explicit in how to setup the problem, so I will make some assumptions to illustrate the effect of scaling.
Given some ''properly scaled'' function $g(x_1, x_2)$ and take $f$ to be the following function
\[
f(x_1,x_2) = g(k_1x_1, k_2x_2)
\]
where $k_1 \gg k_2$ are scale factors to make $f$ poorly scaled.
Given the Hessian has full rank, we can always choose a coordinate system in which the mixed partial derivatives are 0, so we will simply assume that is the case for $\pd[2][][x_1,x_2]$.
This means the Hessian is:
\[
\nabla^2f = \begin{pmatrix}
    k_1^2\pd[2][g][x_1^2]&0\\
    0&k_2^2\pd[2][g][x_2^2]
\end{pmatrix} 
\]
The condition number of a matrix, $\kappa(A)$, is the quotient between its largest and smallest singular values (which for a square matrix is the absolute value of its eigenvalues).
Denoting the set of eigenvalues of a matrix by $\Lambda$, one can write
\[
\kappa(A) = \underset{\lambda_1, \lambda_2 \in \Lambda}{\operatorname{sup}}\left|\frac{\lambda_1}{\lambda_2}\right|
\]
In our case, that becomes:
\[
\kappa \left( \nabla^2 f \right) = \frac{\left|k_1^2\pd[2][g][x_1^2]\right|}{\left|k_2^2\pd[2][g][x_2^2]\right|}
\]
Now if we assume that ''properly scaled'' means that $|\pd[2][g][x_1^2]| \approx |\pd[2][g][x_2^2]|$ at the optimum then that means
\[
\kappa \left( \nabla^2 f \right) \approx \frac{k^2_1}{k_2^2}
\]
So if one axis has a characteristic length of $\sim 10^6$ times larger than the other, the condition number will be $10^{12}$!